\documentclass{situationist-preprint}

\title{Permissivism and Arbitrariness}
\msid{EPI-17-b}
\msfamily{EPI-17}
\msdiscipline{Epistemology}
\mstype{Companion manuscript}
\mskeywords{permissivism; uniqueness; epistemic standards; arbitrariness; immodesty}

\begin{document}
\maketitle

\begin{abstract}
The most developed form of permissivism holds that rational credence is fixed by evidence together with epistemic standards, and that more than one set of standards is rational. We pose a dilemma. Either rationality can evaluate the permitted standards against one another, or it cannot. If it can, one set is better and the permission lapses. If it cannot, the choice between them is arbitrary, and an agent who knows this must expect the credences of the other standards to do as well as her own; with strict propriety, as EPI-17-a shows, that forces the credences to coincide. Arbitrariness of this kind is not the harmless arbitrariness of Buridan's choices.
\end{abstract}

\section{Standards and the dilemma}

Permissivists such as \citet{Schoenfield2014} and \citet{Kelly2014} hold that the rational credence in $p$ is a function of the agent's evidence and her epistemic standards, and that several sets of standards are rational. Two agents with the same evidence and different rational standards may then rationally hold different credences.

\begin{definition}[Rational evaluability]\label{d:eval}
Two sets of standards $S$ and $S'$ are \emph{rationally evaluable} against each other if there is a consideration, available to a rational agent, that favours one as a guide to truth.
\end{definition}

\begin{proposition}[First horn]\label{p:horn1}
If two rational sets of standards are rationally evaluable against each other and one is favoured, a rational agent aware of the consideration does not adopt the other.
\end{proposition}

This is immediate: to adopt the disfavoured standards in the face of an available consideration favouring the others is to ignore a consideration bearing on truth. On the first horn, the plurality of standards does not survive reflection.

\begin{proposition}[Second horn]\label{p:horn2}
If $S$ and $S'$ are both rational and rationally unevaluable against each other, then a rational agent with standards $S$ who knows this is committed to Principle~P2 of EPI-17-a for the credences that $S$ and $S'$ assign on her evidence; and so, by Theorem~3.1 of EPI-17-a, those credences coincide.
\end{proposition}

\begin{proof}[Argument]
Suppose the agent expects the credence $d$ given by $S'$ to be less accurate than her own credence $c$ given by $S$. That expectation is a consideration, available to her, favouring $S$ as a guide to truth on this evidence, contrary to unevaluability. So she does not expect $d$ to do worse: $\mathit{EI}(c,d) \leq \mathit{EI}(c,c)$. This is Principle~P2 of EPI-17-a. By Principle~P1 there (strict propriety), $\mathit{EI}(c,c) < \mathit{EI}(c,d)$ unless $c = d$. Hence $c = d$.
\end{proof}

The permissivist wants unevaluable standards that nonetheless yield different credences. Proposition~\ref{p:horn2} shows that this combination is incoherent for an agent who reflects on it: each rational credence is immodest, and immodesty is a verdict, by the agent's own lights, that her standards are better here. The choice of standards is then arbitrary in a damaging sense, since the agent's own credence treats as worse an option that she has no ground to treat as worse \citep{White2005}.

\section{Objections}

\begin{objection}
Arbitrariness is common and harmless. Buridan's ass must pick one bale arbitrarily, and picking is rational.
\end{objection}

\begin{reply}
In Buridan's case the options are equally good by the agent's lights, so whichever she picks, she does not regard the other as worse. In the permissive case, once a credence is adopted, it regards every alternative as worse in expected accuracy. The symmetry that makes picking harmless is broken by the very credence picked.
\end{reply}

\begin{objection}
The agent may regard her standards as truth-conducive by her own lights and the others' as truth-conducive by theirs; no neutral evaluation is required \citep{Schoenfield2014}.
\end{objection}

\begin{reply}
That concedes the second horn's premise and denies its consequence. If she knows the other standards are rational and that no consideration favours hers, her confidence that hers do better here has no support she can cite. Either she withdraws that confidence, and Principle~P2 of EPI-17-a applies, or she keeps it without ground, which is the arbitrariness the permissivist meant to avoid \citep{Horowitz2014}.
\end{reply}

\section{Conclusion}

Appeal to epistemic standards does not make permissivism stable. Where standards can be compared, comparison eliminates all but one; where they cannot, the agent's credence must endorse the alternatives, and strict propriety then leaves one rational credence for each proposition on each body of evidence.

\begin{thebibliography}{99}\small
\bibitem[Horowitz(2014)]{Horowitz2014} Horowitz, S. (2014). Immoderately rational. \emph{Philosophical Studies}, 167(1), 41--56.
\bibitem[Kelly(2014)]{Kelly2014} Kelly, T. (2014). Evidence can be permissive. In M. Steup, J. Turri and E. Sosa (Eds.), \emph{Contemporary Debates in Epistemology}, 2nd ed. Wiley-Blackwell.
\bibitem[Schoenfield(2014)]{Schoenfield2014} Schoenfield, M. (2014). Permission to believe: Why permissivism is true and what it tells us about irrelevant influences on belief. \emph{No\^us}, 48(2), 193--218.
\bibitem[White(2005)]{White2005} White, R. (2005). Epistemic permissiveness. \emph{Philosophical Perspectives}, 19, 445--459.
\end{thebibliography}

\end{document}
